
\newcommand{\abs}[4]{{#1}\, #2\! : \! #3.\, #4}
\newcommand{\absu}[3]{{#1}\, #2.\, #3}
\mathchardef\mhyph="2D % Define a "math hyphen"

% unicode
\usepackage[utf8]{inputenc}

\DeclareUnicodeCharacter{27F6}{\ensuremath{\longrightarrow}}
\DeclareUnicodeCharacter{25B8}{\ensuremath{\blacktriangleright}}
\DeclareUnicodeCharacter{3BC}{\ensuremath{\mu}}
\DeclareUnicodeCharacter{3BB}{\ensuremath{\lambda}}
\DeclareUnicodeCharacter{3B4}{\ensuremath{\delta}}
\DeclareUnicodeCharacter{3C9}{\ensuremath{\omega}}
\DeclareUnicodeCharacter{3B3}{\ensuremath{\gamma}}
\DeclareUnicodeCharacter{3C4}{\ensuremath{\tau}}
\DeclareUnicodeCharacter{21D2}{\ensuremath{\To}}
\DeclareUnicodeCharacter{22B2}{\ensuremath{\lhd}}
\DeclareUnicodeCharacter{2080}{\ensuremath{_0}}
%\definecolor{darkred}{RGB}{100,10,10}

\newcommand{\arr}[2]{#1 \mathbin{\longrightarrow} #2}

\newcommand{\cata}[1]{\llparenthesis #1 \rrparenthesis}

\newcommand{\utp}[0]{\mathcal{U}}
\newcommand{\lin}[0]{\lambda^{\iota\nu}}
\newcommand{\interp}[1]{\llbracket #1 \rrbracket}
\newcommand{\interpl}[0]{\llbracket}
\newcommand{\interpr}[0]{\rrbracket}
\newcommand{\choice}[0]{\xi}
\newcommand{\simto}[0]{\triangleright}
\newcommand{\elcap}[0]{\cap}

\newcommand{\dedu}[0]{\triangleright}
\newcommand{\ctxtup}[1]{\ulcorner #1 \urcorner}
\newcommand{\ctxtdn}[1]{\llcorner #1 \lrcorner}
\newcommand{\ctxtmv}[1]{\langle #1 \rangle}

\newcommand{\pcase}[1]{\noindent\underline{#1:}}

\newcommand{\adbmal}[0]{\reflectbox{$\lambda$}}

\newcommand{\casetri}[0]{\RHD}
\newcommand{\lke}[0]{\lhd}

\newcommand{\lam}[2]{\lambda\, #1.\, #2}
\newcommand{\mal}[2]{\abdmal\, #1.\, #2}
\newcommand{\all}[2]{\forall\, #1.\, #2}
\newcommand{\Lam}[2]{\Lambda\, #1.\, #2}
\newcommand{\thereis}[2]{\exists\, #1.\, #2}
\newcommand{\nuu}[2]{\nu\, #1.\,#2}
\newcommand{\inn}[2]{\langle #1 \rangle_{#2}}
\newcommand{\exel}[2]{\adbmal\,#1.\,#2}
\newcommand{\recc}[2]{\textit{rec}\, #1.\, #2}
\newcommand{\mew}[2]{\mu_{#1}\, #2}
\newcommand{\convl}[0]{\blacktriangleleft}
\newcommand{\convr}[0]{\blacktriangleright}
\newcommand{\conjin}[0]{\smalltriangleleft}
\newcommand{\conjout}[0]{\smalltriangleright}
\newcommand{\dcdot}[0]{\mathbin{\ast}}
\newcommand{\To}[0]{\Rightarrow}
\newcommand{\releq}[0]{\topdoteq}

\newcommand{\omga}[4]{\omega\ #1(#2)\,:\,#3.\ #4}
\newcommand{\clause}[2]{#1 \to #2}

\newcommand{\hysep}[0]{\mathbin{\blacklozenge}}

% \newcommand{\todo}[1]{\textbf{[#1]}}
\newcommand{\rref}[1]{\{\ref{#1}\}}

%\newtheorem{definition}[theorem]{Definition}

% \newenvironment{mycomment}
%     {\color{blue} \it
%     }
%     {
%     }
\newcommand{\shift}[2]{{\uparrow^{#1}_{#2}}}
\newcommand{\textblue}[1]{\textcolor{blue}{#1}}

\newcommand{\lcata}[0]{\ensuremath{\mathcal{L}_{\textit{cata}}}}


\newcommand{\gsep}{\mspace{8mu}\lvert\mspace{8mu}}

% types
\newcommand{\sBool}{\textit{Bool}}
\newcommand{\sNat}{\textit{Nat}}
\newcommand{\sNatList}{\textit{NatList}}
\newcommand{\sAlg}{\Rightarrow}
\newcommand{\sArr}{\to}
\newcommand{\sFunctorApp}[2]{#1 ~ #2} % spaced
% \newcommand{\sFunctorApp}[2]{#1(#2)} % parenthesized
\newcommand{\sMu}[1]{\mu #1}
\newcommand{\sFunctorConSig}[3]{(#1) \sArr \sFunctorApp{#2}{#3}}
\newcommand{\sCoerceSig}[2]{#1 <:: #2}
% terms
\newcommand{\sLet}[3]{\text{let}~ #1 = #2 ~\text{in}~ #3}
% \newcommand{\sCoerce}[2]{\text{coerce}~ #1 ~ #2}
\newcommand{\sCoerce}[2]{#1 - #2}
% coercions
\newcommand{\sCata}[2]{\text{cata}~ #1 ~ #2}
\newcommand{\sRoll}[1]{\text{roll}~ #1}
\newcommand{\sUnroll}[1]{\text{unroll}~ #1}
\newcommand{\sSubCov}[2]{\text{cov}[#1] ~ #2}
\newcommand{\sSubData}[1]{\text{id}[\sMu{#1}]}
\newcommand{\sSubArr}[2]{\text{arr}~ #1 ~ #2}
\newcommand{\sSubAlg}[2]{\text{alg}[#1] ~ #2}
\newcommand{\sSubRecVar}[1]{\text{id}[#1]}
% constructors
\newcommand{\sTrue}{\textit{True}}
\newcommand{\sFalse}{\textit{False}}
\newcommand{\sZero}{\textit{Zero}}
\newcommand{\sSuc}{\textit{Suc}}
\newcommand{\sNil}{\textit{Nil}}
\newcommand{\sCons}{\textit{Cons}}

% helpers
\newcommand{\freshMetavariables}[1]{#1 ~\text{fresh}}
\newcommand{\isDatatype}[1]{#1 ~\text{data}}
\newcommand{\atPos}[1]{^{#1}}
\newcommand{\constructorOf}[4]{(#1 : \sFunctorConSig{#2}{#3}{#4}) \in \mathscr{K}(#3, #4)}
% \newcommand{\atPos}[1]{}

\newcommand{\Match}{\textbf{match }}
\newcommand{\With}{\textbf{ with}}
\newcommand{\Case}{\textbf{case }}
\newcommand{\Func}[1]{\textsc{#1}}

\newcommand{\todo}[1]{{{\contour{red}{\textbf{\color{black}TODO}}} \color{red}#1}}

% \newcommand{\todoHenry}[1]{%
%   \begin{tcolorbox}[
%       colback=olive!15,     % Light olive background
%       colframe=olive!50,    % Muted olive border
%       arc=0mm,            % Sharp corners
%       boxrule=0.5pt,      % Border thickness
%       left=6pt, right=6pt, top=6pt, bottom=6pt, % Inner padding
%       before=\par\smallskip, after=\par\smallskip % Paragraph-like spacing
%     ]
%     \todo{Henry}
%     \\[1em]
%     #1
%   \end{tcolorbox}%
% }
\newcommand{\todoHenry}[1]{%
  \todo{Henry}: #1%
}

% \newcommand{\todoAaron}[1]{%
%   \begin{tcolorbox}[
%       colback=blue!15,     % Light blue background
%       colframe=blue!50,    % Muted blue border
%       arc=0mm,            % Sharp corners
%       boxrule=0.5pt,      % Border thickness
%       left=6pt, right=6pt, top=6pt, bottom=6pt, % Inner padding
%       before=\par\smallskip, after=\par\smallskip % Paragraph-like spacing
%     ]
%     \todo{Aaron}
%     \\[1em]
%     #1
%   \end{tcolorbox}%
% }
\newcommand{\todoAaron}[1]{%
  \todo{Aaron}: #1%
}
